\section{Where the hypotheses fail in eight dimensions: the counterexamples}
\label{sec:catalogue}

Table~\ref{tab:catalogue} collects the constructions of Secs.~\ref{sec:octonions} to \ref{sec:susy}
and of the superfluid of Sec.~\ref{sec:superfluid}. Each one combines Lorentz and internal symmetry
in a way that the conclusion of the no-go theorems forbids, and each one violates at least one
hypothesis of Table~\ref{tab:hypotheses}; the last column names the machine checks that establish it.
None is a counterexample to a theorem whose hypotheses hold---there is none---but together they map
out how much room eight dimensions, and the octonions in particular, leave.

Three lessons stand out. First, the specifically octonionic mechanisms act on the \emph{first}
hypothesis, H0: non-associativity removes octonionic charges from operator algebras altogether
(Theorem~1), and triality, an automorphism of the Lie algebra rather than an operator, exists only
for the split signature $4+4$, which violates H0 (indefinite metric) and H1 (four times) anyway.
Second, on $7+1$ the hypotheses H0 and H2 to H5 hold for any reasonable field theory with a mass gap,
and the theorem of Pelc and Horwitz applies: every combination of Lorentz and internal symmetry on
$7+1$ requires either odd generators (supersymmetry) or a background that breaks the eight-dimensional
Poincar\'e group (a combined manifold, a medium, a condensate). Third, the combined manifolds turn a
part of the eight-dimensional Lorentz group into an internal group of the four-dimensional world; that
internal group can then be combined freely with other internal symmetries, and the superfluid of
Sec.~\ref{sec:superfluid} does so: its chiral phase leaves unbroken a sum of an eight-dimensional
Lorentz rotation and the fermion number, in a vacuum that is invariant under the four-dimensional
Lorentz group.

\begin{table*}[t]
\caption{The constructions of this paper that combine Lorentz and internal symmetry, the hypothesis
of Table~\ref{tab:hypotheses} that each violates, and the checks of Appendix~\ref{app:checks} that
establish it.}
\label{tab:catalogue}
\begin{ruledtabular}
\begin{tabular}{>{\raggedright\arraybackslash}p{0.27\textwidth}l>{\raggedright\arraybackslash}p{0.33\textwidth}>{\raggedright\arraybackslash}p{0.12\textwidth}l}
construction & spacetime & what is combined & failing hypothesis & checks\\
\hline
octonionic charges (Malcev algebra) & both & imaginary units as charges; their associative envelope is $\mathfrak{so}(8)$ or $\mathfrak{so}(4,4)$, the rotations themselves & H0 (no Jacobi identity) & O-7 to O-9, O-12\\
triality & $4+4$ & vectors exchanged with half-spinors by an outer automorphism & H0, H1 & T-1 to T-6\\
no triality & $7+1$ & (proof that the order-three automorphism does not exist) & --- & T-7a to T-7d\\
exchange of world and fibre & $4+4$ & left multiplication by a split unit swaps the Lorentz algebras of world and fibre & H0, H1 & K-6\\
non-compact fibre algebra & $4+4$ & internal $\mathfrak{so}(1,3)$ from the Lorentz group, non-unitary & H0 & K-2, K-7\\
Kaluza-Klein tachyons & $4+4$ & internal momenta lower the mass below zero & H2 & K-9\\
reduction to world times fibre & $7+1$ & internal $\mathfrak{so}(4)$ from the Lorentz group & H1 (in eight dimensions) & K-1, K-3\\
mass splitting & $7+1$ & a Lorentz rotation changes four-dimensional masses & H7 & K-8\\
$\mathrm G_2$ locking & $7+1$ & space rotations locked to fibre rotations by the octonionic three-form & H1 (a medium) & K-5a to K-5c\\
super-Poincar\'e algebra with R-symmetry & $7+1$ & R-charge rotates the supercharges, whose anticommutator is a momentum & H0 (odd generators) & S-1, S-2a, S-4\\
super-Poincar\'e algebra & $4+4$ & as above, without positive energy & H8, H1 & S-2a, S-4\\
tensorial charges & $7+1$ & charges commuting with all momenta but not Lorentz scalars & H6 (abelian) & S-2b, S-3\\
internal-vector condensate & both & stabilizer mixes world and fibre directions in eight dimensions & H1 (in eight dimensions) & K-10, P-10\\
chiral spin-zero superfluid & $7+1$ & a fibre rotation (part of the eight-dimensional Lorentz group) locked to the fermion number & H1 (in eight dimensions) & P-3a, P-10\\
\end{tabular}
\end{ruledtabular}
\end{table*}
